\section{Introduction}
\label{sec:intro}

\subsection{Clinical motivation}
\label{sec:intro:clinical}

Spectral Doppler ultrasound infers blood-flow velocity from the frequency shift
$\fd$ that moving red blood cells impose on the insonating beam. The shift and
the velocity are tied by the Doppler equation,
%
\begin{equation}
  \fd \;=\; \frac{2\,f_0\,v\,\cos\theta}{c},
  \label{eq:doppler}
\end{equation}
%
where $f_0$ is the transmit frequency, $v$ the flow speed, $c$ the speed of
sound in soft tissue, and $\theta$ the \emph{Doppler angle} between the beam and
the local direction of flow. Rearranging for the clinically reported quantity
gives $v = c\,\fd / (2 f_0 \cos\theta)$, so every velocity estimate carries a
multiplicative $1/\cos\theta$ correction.

The machine does not measure $\theta$. The sonographer sets it by hand, aligning
a steering cursor with the vessel wall on the live B-mode image. This manual step
makes $\theta$ an operator-dependent source of error, and the $1/\cos\theta$
factor amplifies that error nonlinearly. At the steep angles routinely used in
carotid scanning the sensitivity $\mathrm{d}v/\mathrm{d}\theta \propto \tan\theta$
grows without bound, so a few-degree cursor mistake near
$\theta\!\to\!\SI{90}{\degree}$ translates into a large velocity error. Carotid
peak-velocity thresholds feed directly into stenosis grading, so reducing the
operator dependence of the angle measurement is clinically relevant.

The motivation is not hypothetical. Improper angle correction is repeatedly cited
among the most common Doppler errors, and as many as \SI{35}{\percent} of vascular
laboratories applying for accreditation have been flagged for improper
angle-correction technique~\citep{Saad2008}, a recurring cause of delayed
accreditation decisions. We note that an image-only angle reader automates the
operator's cursor placement; it does not establish true flow direction
independently of the reference reading, a point we return to in
\S\ref{sec:discussion}.

\subsection{Reading the angle from the image}
\label{sec:intro:idea}

The angle a sonographer sets is, by construction, the geometric relationship
between the beam axis and the vessel as it appears on the grayscale B-mode frame.
That relationship is present in the image itself, which suggests estimating
$\theta$ directly from a single B-mode image, with no color Doppler overlay and
no explicit vessel segmentation. Cast this way, angle correction becomes a
supervised regression problem, and the manual cursor placement is replaced by an
automated reading that is reproducible across repeats of the same image. The
ground truth is still the operator-set angle, so the model can be no more accurate
than its reference labels. The practical difficulty is that clinical carotid
datasets are small, and learning a continuous geometric quantity from a few dozen
images is the regime where modern deep networks tend to overfit.

\subsection{Prior work and the present contribution}
\label{sec:intro:prior}

Automating the Doppler angle is a long-standing goal. Earlier approaches operated
on \emph{color} flow images with an explicit vessel geometry. Hirsch et al.\
estimated the vessel orientation by a multi-scale principal-component analysis to
place the Doppler gate \citep{Hirsch2006}. Saad et al.\ segmented the vessel in
color-Doppler images and derived the angle from a skeleton of the segmented lumen,
reporting strong agreement with manual readings \citep{Saad2008}. Both rely on
color information and a segmentation step, so artifacts in the color-flow image
propagate into the angle estimate. Deep learning has meanwhile been adopted across
ultrasound, including fetal standard-plane localization \citep{Chen2015},
echocardiographic view classification \citep{Madani2018}, and sensorless 3D
freehand reconstruction \citep{Prevost2018}. This suggests that the geometric cues
a sonographer uses can be learned directly from the image.

Building on this, Patil and Anand introduced an image-only framing and showed that
a deep convolutional network could read $\theta$ directly from a \emph{grayscale}
B-mode carotid image, with no color Doppler and no segmentation, removing the
manual cursor step \citep{Patil2019}. Because such an estimator uses only the
formed B-mode image and no knowledge of the vendor-specific signal-processing
path, it does not depend on a particular color-flow pipeline. We emphasize that
this is an argument about inputs, not a demonstration of cross-device transfer:
the data here come from a single source, so any claim of device- or
vendor-agnostic deployment would have to be tested separately. To cope with data
scarcity the original study began from a base set of $84$ longitudinal
common-carotid B-mode images acquired from roughly ten volunteers and enlarged it
by rotating each image over $[-60,60]\si{\degree}$ in \SI{5}{\degree} increments.
This produces $25$ oriented views per image and a working corpus of about
$2{,}100$ examples. The known rotation supplies an \emph{exact} angle label for
every view, so the augmentation is self-labelling, which is rare for a regression
target. On this corpus a frozen ImageNet backbone \citep{Deng2009} with a
lightweight regression head, used purely as a fixed feature extractor in the
standard transfer-learning manner, reached single-digit errors; the best reported
configuration achieved roughly \SI{2.87}{\degree} mean absolute error (MAE) and
\SI{4.03}{\percent} mean absolute percentage error (MAPE).

This paper treats Doppler-angle estimation as an engineering target and asks how
far a reproducible, frozen-transfer estimator can be pushed on the same data. We
reproduce the original pipeline and identify the design choice that makes it work,
tune and ensemble the estimator under two sampling protocols, run an architecture
comparison, and evaluate the estimator with calibrated uncertainty and a classical
baseline.

There is no single correct way to score accuracy on a small, heavily augmented
cohort, so we report two sampling protocols. \emph{Image-level sampling} is the
original study's protocol: the $\sim\!2{,}100$ rotated views are partitioned into
train and test at the image level, so the resulting error measures accuracy across
the population of orientations and imaging conditions the augmentation spans.
\emph{Patient-level sampling} holds out whole volunteers (grouped
cross-validation), so it measures generalization to a previously unseen patient.
The two answer different questions, and the patient-level number is higher because
cross-subject transfer on $\sim\!10$ volunteers is intrinsically harder. We tune
each protocol to its own best and report both. With only about ten volunteers,
grouped five-fold cross-validation holds out roughly two volunteers per fold, so
the patient-level estimates carry large fold-to-fold variance and the point
estimates below should be read accordingly. We report the per-fold spread in
\S\ref{sec:results}. The image-to-patient gap quantifies how much of the
within-population accuracy is anatomy-specific
(\S\ref{sec:methods}, \S\ref{sec:discussion}).

\paragraph{Faithful, reproducible replication.}
We first reproduce the original protocol end to end on a modern, fully scripted
stack (Keras\,3 / JAX, fixed seed). A frozen DenseNet201 backbone
\citep{Huang2017} with an orientation-preserving \emph{grid-pooling} head reaches
\SI{5.84}{\percent} MAPE (\SI{3.77}{\degree} MAE, $R^2=0.982$) under image-level
sampling, on the original single augmented split and without any fine-tuning
(Table~\ref{tab:replication}).

\paragraph{The pooling layer.}
That replication hinges on the pooling layer. Conventional global average pooling
(GAP) averages the convolutional feature map over all spatial positions, which is
approximately rotation-invariant. That is the wrong inductive bias when the target
\emph{is} an orientation. Grid pooling instead average-pools the feature map to a
small $G\times G$ grid and flattens it, preserving coarse spatial layout. On the
same frozen DenseNet201 this roughly halves the error under image-level five-fold
cross-validation, from \SI{10.85}{\percent} to \SI{4.58}{\percent} MAPE; the
original single augmented split, scored separately, reaches \SI{5.84}{\percent}
MAPE (\SI{3.77}{\degree} MAE) (Table~\ref{tab:replication}). The two figures use
different aggregation schemes (a cross-validation mean versus a single held-out
split) and are not directly comparable. The reading is the same in both: the
frozen backbone already encodes vessel orientation, and the gain comes from not
pooling it away.

\paragraph{Per-protocol tuning and ensembling.}
We tune the regression head (depth, width, dropout, $L_2$, batch normalization,
learning rate, batch size, patience) with the Optuna TPE sampler
\citep{Akiba2019,Bergstra2011}, scoring each protocol on its own five-fold
cross-validation over cached frozen features. One feature extraction per backbone
serves both protocols. We then combine five tuned backbones into mean and stacked
\citep{Wolpert1992} ensembles over out-of-fold predictions
(Table~\ref{tab:dual_protocol}). The stacked ensemble is the strongest estimator
under both protocols: \SI{2.79}{\percent} MAPE (\SI{1.96}{\degree} MAE,
$R^2=0.995$) image-level and \SI{8.53}{\percent} MAPE (\SI{5.93}{\degree} MAE,
$R^2=0.952$) patient-level. These ensemble figures pool out-of-fold predictions
before scoring, rather than averaging per-fold scores, so they are not a strict
cross-validation mean and stacking on out-of-fold predictions can be mildly
optimistic; we keep the distinction in view when comparing them to the replication
numbers. A single tuned DenseNet201 reaches \SI{3.00}{\degree} MAE and
\SI{4.03}{\percent} MAPE, close to the prior best of \SI{2.87}{\degree} MAE and
slightly worse on MAE. The ensemble improves on this and also reports the
previously unmeasured patient-level number.

\paragraph{Architecture comparison.}
Larger and newer encoders are often assumed to be better, so we compare frozen
classic and modern ImageNet backbones under patient-level five-fold
cross-validation (Fig.~\ref{fig:bakeoff}). On the extraction-matched comparison,
frozen DenseNet201 \citep{Huang2017} leads at \SI{14.1}{\percent} MAPE, ahead of
the best modern encoder, ConvNeXt-Base \citep{Liu2022} (\SI{15.7}{\percent}), with
EfficientNet/EfficientNetV2 \citep{Tan2019} trailing at
\SIrange{17}{21}{\percent}. With only $84$ base images, no newer or larger backbone
improves on a frozen 2017 DenseNet201, which we carry forward.

\paragraph{Clinical-grade evaluation.}
Finally we report uncertainty and agreement in clinically legible terms.
Split-conformal prediction \citep{Vovk2005,Lei2018} on a patient-disjoint
calibration split yields \SI{90}{\percent} intervals of $\pm\SI{20.50}{\degree}$,
with empirical coverage of \SI{95.2}{\percent} on the held-out test set. Conformal
prediction guarantees marginal coverage in expectation over calibration draws, not
the realized coverage on one split. With only $\sim\!10$ volunteers a single split
gives a coverage estimate with wide binomial uncertainty, and the realized
\SI{95.2}{\percent} sits above the \SI{90}{\percent} target, so the intervals are
on the conservative side here.

A Bland--Altman analysis \citep{BlandAltman1986} reports a $-\SI{4.31}{\degree}$
bias (method minus reference: the model reads slightly low) with limits of
agreement $[-24.25,\,+15.63]\si{\degree}$. This compares the method against the
\emph{single} reference reading available per image, so it is method-vs-reference
agreement rather than an inter-observer study, and the limits combine model error
with reference noise that we cannot separately estimate. The bias is directional,
which would mean a systematic rather than random velocity offset; we discuss its
clinical implication in \S\ref{sec:discussion}.

The headline image-level numbers are computed over the rotation-augmented
population. On real unrotated base images the MAE is \SI{7.80}{\degree}, and
rotation test-time augmentation nearly halves it to \SI{4.72}{\degree} MAE
(circular-median over rotations). We foreground this base-image figure as the most
clinically relevant per-image performance. Grad-CAM \citep{Selvaraju2017} is
consistent with the network attending to the vessel wall. As a check on the
learned representation, a classical, image-only structure-tensor angle prior
reaches \SI{3.16}{\degree} MAE on the narrow base-angle band, and a circular
($2\theta$) fusion of the learned and classical estimates reaches
\SI{2.72}{\degree} MAE, improving on either alone on this point estimate. We read
the fusion gain as suggestive of partly complementary cues, pending the variance
analysis in \S\ref{sec:results}.

\subsection{Contributions}
\label{sec:intro:contrib}

In summary, this paper makes the following contributions.
%
\begin{itemize}[leftmargin=1.4em]
  \item \textbf{A faithful, fully scripted replication.} We reproduce the prior
  pipeline (\SI{5.84}{\percent} MAPE, frozen DenseNet201) and isolate the design
  choice that makes it work: an orientation-preserving grid-pooling head in place
  of global average pooling, so that a frozen backbone retains the orientation
  signal the target depends on.
  \item \textbf{The strongest estimator we measured, under two protocols.} Via
  per-protocol Optuna head tuning and a five-model ensemble we report
  \SI{2.79}{\percent} MAPE (\SI{1.96}{\degree} MAE) under image-level sampling and
  \SI{8.53}{\percent} MAPE (\SI{5.93}{\degree} MAE) under the stricter
  patient-level sampling, each tuned to its own best, alongside the previously
  unmeasured cross-subject number.
  \item \textbf{An architecture comparison.} A frozen DenseNet201 outperforms newer
  ConvNeXt and EfficientNet/EfficientNetV2 backbones, indicating that newer is not
  necessarily better for small-data frozen transfer on this task.
  \item \textbf{A calibrated clinical-grade evaluation.} We provide conformal
  prediction intervals, a method-vs-reference Bland--Altman analysis stated with
  its single-reference caveat, rotation test-time augmentation, and Grad-CAM
  attribution, plus a classical structure-tensor prior whose circular fusion with
  the learned model reaches \SI{2.72}{\degree} MAE.
\end{itemize}
